\subsection{関数の平均値}
$a<b$ のとき,
\[
 f_{\mathrm{av}}=\frac1{b-a}\int_a^b f(x)\dd x
\]
を $f$ の区間 $[a,b]$ における\textbf{平均値}という.等分割のリーマン和を区間の長さで割ると $\frac1n\sum_{k=1}^{n}f(\xi_k)$ になるので,有限個の値の平均を細かくしていったものと解釈できる.微分で用いた平均変化率とは別の量である.

$f$ の最小値を $m$,最大値を $M$ とすると,積分の評価から $m\le f_{\mathrm{av}}\le M$ である.連続関数は最小値と最大値の間の値をすべてとるから,ある $c\in[a,b]$ で
\[
 \int_a^b f(x)\dd x=f(c)(b-a)
\]
が成り立つ.これを\textbf{積分の平均値の定理}という.高さを一定の $f(c)$ にした長方形と,元の関数の定積分が等しくなるという意味である.



